\chapter{Injection-Time Thinning and the Opacity Mechanism}
\label{ch:thin}

\section{The crowding problem}
\label{sec:thin:crowding}

A two-view prediction covers each visible surface once. A SLAM map covers it
once \emph{per keyframe that saw it}, and on a typical indoor sequence that is
tens of times. Nothing in the pipeline reconciles these copies: each is
predicted independently, placed by its own keyframe's pose, and injected with
whatever opacity its forward pass produced.

Under \cref{eq:alphacomp} the consequence is specific rather than generic. If
the copies agreed in colour and depth, stacking would still change effective
opacity: $K$ copies with opacity $\alpha$ have combined opacity
$1-(1-\alpha)^K$. Inconsistent copies additionally differ in depth by the per-forward scale gauge
(\cref{sec:neg:frontend} measures this at ${\sim}5\%$), in colour by the
source-frame photometric anchoring of \cref{eq:colorresidual}, and in extent by
the head's own scale prediction. A ray therefore accumulates opacity from
several slightly-displaced, slightly-different surfaces before its
transmittance runs out, and what should be one sharp surface composites into
haze.

The released checkpoint makes this worse in a way that is visible in its
statistics rather than in its images: its predicted opacity is
\emph{saturated}, with median $1.000$ and every Gaussian above $0.9$
(\cref{sec:adapt:aug}). A saturated opacity channel cannot express ``this
copy should contribute less''. The lever below supplies that expression from
outside the network.

\section{Insertion-time opacity attenuation}
\label{sec:thin:lever}
\label{sec:method:thin}

A SLAM map covers each surface with predictions from many keyframes. Under
\cref{eq:alphacomp} with a black background, a surface covered by many
over-solid Gaussians composites to haze rather than to a sharper surface. We
therefore fade opacity at the moment a Gaussian enters the map, in proportion
to how little the pointmap confidence vouches for it.

Confidences are not calibrated across keyframes, so we rank-normalise
\emph{within} each keyframe. For a keyframe with $N$ Gaussians and confidence
values $\{c_n\}$, let $r_n$ be the rank of $c_n$ and
$\hat{c}_n = r_n/(N-1) \in [0,1]$. The injected opacity is
\begin{equation}
\alpha'_n \;=\; \alpha_n \cdot
\mathrm{clip}\Bigl(1 - 2\lambda\,(1-\hat{c}_n),\; 0.1,\; 1\Bigr),
\label{eq:thin}
\end{equation}
with dose $\lambda{=}0.45$ in the thinning experiments. External comparisons
disable thinning with $\lambda=0$. The lower clip at $0.1$ retains every
primitive and its refinement gradient.

\subsection{Relation to pairwise training}
An opacity penalty added to \cref{eq:splatt3rloss} would act during pair
training. Crowding appears later, after SLAM accumulates
keyframes over a surface, and \cref{eq:splatt3rloss} never sees more than two
context views. Consistent with this, we measured that the refiner
un-saturates opacity on its own even with the released head (median
$1.0\rightarrow0.26$). Attenuation therefore starts the map closer to an
opacity state that the optimiser can reach. A dose--response check at two
thinning depths (opacity knees $0.9$ and $0.6$) gave $-21.0\%$ and
$-15.8\%$ LPIPS. Benefits persist for these heads, supporting complementarity
in this experiment. A different training objective or sequence-level
training procedure might still learn an equivalent adjustment.

\section{Opacity, colour, and visibility}
\label{sec:method:mech}

The DC colour is an SH residual added to the encoding of input RGB, with
clamping at conversion (\cref{eq:colorresidual}). Both the residual and
opacity can respond to photometric perturbations. Opacity additionally
controls which surfaces remain visible along a ray.

For an isolated primitive over black, greater opacity increases its
contribution. In a multicolour stack it may instead obscure brighter content.
The background therefore does not imply a monotone relation between opacity
and brightness. We test their correlation without assigning a unique cause.

We registered a sign-specific prediction. Let $\ell$ be input
luminance and $\rho$ the Spearman correlation between $\alpha$ and $\ell$ over
predicted Gaussians. A head trained with independently darkened inputs
(synthetic vignetting and white-balance jitter) was predicted to have
$\rho<0$, whereas a real-sensor head without synthetic darkening was
predicted to have $\rho>0$. \Cref{sec:exp:mech} reports both signs in the
five tested heads: two Replica augmentation heads and three real-data
heads. These are five models, not five independent dataset families.
Correlations are consistent with opacity responding to photometry, but do
not identify a unique causal mechanism.

\section{Controls}
\label{sec:thin:ablations}

\subsection{Confidence weighting}

\Cref{eq:thin} fades in proportion to rank-normalised confidence, and the name
``confidence-weighted thinning'' invites the reading that the confidence signal
is what makes it work. We measured the alternative --- the same total dose
applied \emph{uniformly}, with no confidence term at all --- and against
confidence weighting it is statistically a coin flip: $n{=}12$ online cells,
median difference $-0.06\%$.

Reducing injected opacity is the robust effect in this study. The
confidence-weighted form concentrates reduction where the pointmap is least
sure, but the uniform control performs similarly.

\subsection{Shape-weighted attenuation}

A second opacity lever reduces each Gaussian's opacity in proportion to how
elongated it is relative to its local surface sampling. It targets trailing
primitives near depth discontinuities. In measurement it fades $64$--$98\%$
of all Gaussians by $39$--$65\%$ of their opacity. Its effect is therefore
broad rather than sparse.

LPIPS improves on the three tested sequences by $-12.7\%$, $-6.6\%$, and
$-1.3\%$, with a PSNR change of $-0.04$ to $-0.22$\dB{}.
The LPIPS gain is largest in scenes with more streaked geometry.
